% =============================================================
% setup.tex - Dokumentklasse, Pakete & Schriftarten
% Ersetzt: documentclass.tex + load_packages.tex + Schriftart-settings.tex
% =============================================================

% --- Warnungen unterdrücken ---
\RequirePackage{silence}
\WarningFilter*{titlesec}{Non standard sectioning command detected}
\WarningFilter*{titlesec}{Usage of package `titlesec'}
\WarningFilter*{titlesec}{Activating an ugly workaround}

% =============================================================
% 1. DOKUMENTKLASSE
% =============================================================

\documentclass[
  12pt,
  a4paper,
  twoside,
  openright,
  bibliography=totoc,
  listof=totoc,
  captions=tableheading,
  headinclude,
  numbers=noenddot,
  titlepage=true,
  abstract=true,
  version=last,
  final
]{scrreprt}

\KOMAoptions{
  parskip=half,
  chapterprefix=true,
  headsepline=true,
  footsepline=false,
  DIV=12,
  BCOR=0mm,
  headings=normal,
  footnotes=multiple
}

% Überschriften-Stile
\addtokomafont{disposition}{\rmfamily\bfseries}
\addtokomafont{chapter}{\large}
\addtokomafont{section}{\normalsize}
\addtokomafont{subsection}{\small}
\addtokomafont{caption}{\small\itshape}
\addtokomafont{captionlabel}{\bfseries}

% Nummerierungstiefe
\setcounter{secnumdepth}{3}
\setcounter{tocdepth}{3}

% Kapitel-Format
\renewcommand*{\chapterformat}{%
  \chapappifchapterprefix{\nobreakspace}\thechapter\autodot\enskip
}

% LuaTeX-Prüfung
\usepackage{iftex}
\ifLuaTeX
  \RequirePackage{luatex85}
\else
  \PackageError{setup}{Diese Vorlage erfordert LuaLaTeX!}{}
\fi

% =============================================================
% 2. SCHRIFTARTEN
% =============================================================

\usepackage{fontspec}

% Hauptschrift: Bookerly aus fonts/ (liegt im Repo, muss nicht installiert sein)
\IfFileExists{fonts/Bookerly-Regular.ttf}{
  \setmainfont{Bookerly}[
    Path = fonts/,
    Extension = .ttf,
    UprightFont = Bookerly-Regular,
    BoldFont = Bookerly-Bold,
    ItalicFont = Bookerly-RegularItalic,
    BoldItalicFont = Bookerly-BoldItalic
  ]
}{
  \setmainfont{TeX Gyre Heros}[
    Scale=1.0,
    Ligatures=TeX,
    UprightFont=*-Regular,
    BoldFont=*-Bold,
    ItalicFont=*-Italic,
    BoldItalicFont=*-BoldItalic,
  ]
}

% Monospace-Schrift für Code
\IfFontExistsTF{Consolas}{
  \setmonofont{Consolas}[Scale=0.9]
}{
  \IfFontExistsTF{DejaVu Sans Mono}{
    \setmonofont{DejaVu Sans Mono}[Scale=0.9]
  }{
    \setmonofont{TeX Gyre Cursor}[Scale=0.9]
  }
}

% Typografie
\usepackage{microtype}
\microtypesetup{protrusion=true, expansion=true}

% Trennverhalten
\tolerance=2000
\emergencystretch=1em
\hyphenpenalty=5000
\clubpenalty=10000
\widowpenalty=10000

% =============================================================
% 3. PAKETE
% =============================================================

% Sprache
\usepackage{polyglossia}
\setmainlanguage[spelling=new]{german}
\setotherlanguage{english}

% Kommentare
\usepackage{comment}

% Mathematik
\usepackage{amsmath, amssymb, amsfonts, amsthm, amsbsy}
\usepackage{mathtools}
\usepackage{unicode-math}

% Tabellen
\usepackage{booktabs}
\usepackage{multirow}
\usepackage{array, longtable}
\usepackage{tabularray}
\usepackage{lscape}
\usepackage{pdflscape}

% Grafiken
\usepackage{graphicx}
\usepackage{qrcode}
\graphicspath{{media/}}
\usepackage{float}
\usepackage{placeins}
\usepackage{rotating}
\usepackage{subcaption}
\usepackage{wrapfig}
\usepackage{adjustbox}

% Symbole & TikZ
\usepackage{fontawesome5}
\usepackage{tikz}
\usetikzlibrary{positioning, shapes.geometric, arrows.meta}

% Text & Layout
\usepackage{ragged2e}
\usepackage{setspace}
\setstretch{1.1}
\usepackage{enumitem}
\usepackage{nowidow}

% Bibliographie
\usepackage[
  backend=biber,
  style=numeric,
  sorting=none,
  sortcites=true,
  giveninits=true,
  maxbibnames=99
]{biblatex}

% Glossar
\usepackage[acronym,toc,nogroupskip]{glossaries}

% PDF & externe Inhalte
\usepackage{pdfpages}
\usepackage{todonotes}

% URLs & Links
\usepackage{url}
\usepackage{xurl}

% Beschriftungen
\usepackage[
  format=plain,
  labelfont=bf,
  justification=centering,
  tableposition=top,
  figureposition=bottom,
  font=small
]{caption}

% Farben
\usepackage{xcolor}

% Quotes & Platzhalter
\usepackage{csquotes}
\MakeOuterQuote{"}
\usepackage{lipsum}

% Theorem-Definitionen
\theoremstyle{definition}
\newtheorem{definition}{Definition}[chapter]
\newtheorem{beispiel}[definition]{Beispiel}
\newtheorem{lemma}[definition]{Lemma}
\newtheorem{satz}[definition]{Satz}

% Utility-Befehle
\usepackage{xspace}
\newcommand{\zB}{z.\,B.\xspace}
\newcommand{\dhh}{d.\,h.\xspace}
\newcommand{\ua}{u.\,a.\xspace}

% Hilfsbefehle
\newcommand{\tick}{\textcolor{green!60!black}{✓}}
\newcommand{\cross}{\textcolor{red!60!black}{✗}}

% Rotierte Seite
\newenvironment{rotatepage}{%
  \clearpage\begin{landscape}
}{%
  \end{landscape}\clearpage
}

% LuaTeX-spezifische Pakete
\usepackage{luatex85}
\ifLuaTeX
  \usepackage{luacode}
  \usepackage{shellesc}
\fi
